(1) 一维探针在晶格周期势 \(V_0[1-\cos(kx)]\) 中运动，\(k=2\pi/a\)，又以劲度 \(K>0\) 的弹簧接到匀速驱动端 \(X(t)=vt\)，\(v>0\)。总势 \(V(x,X)=V_0[1-\cos(kx)]+K(x-X)^2/2\)。求平衡方程、局部稳定条件和消除准静态黏滑多稳态的临界刚度。
(2) 在 \(0<K<K_c\) 下找出随 \(X\) 增大时某势阱的前向失稳（折叠）点 \(x_s,X_s\)，其中 \(\sin(kx_s)>0\)，并给出其适用的相位范围。
(3) 设 \(\delta X=X_s-X>0\) 很小，展开总力至 \(x_s\) 的二次项，求稳定点与鞍点的距离、趋零的势阱曲率，以及两者之间势垒随 \(\delta X\) 的幂律及系数。
(4) 加入过阻尼运动 \(\zeta\dot x=-\partial_xV\)，由非自治势能严格推导外驱功率—储能变化—耗散关系。解释在零温慢驱动下何时发生跳跃，以及在什么情况下惯性、热激活会破坏该理想零温准静态公式。
